\documentclass[10pt,a4paper]{amsart}
\usepackage[margin=19mm]{geometry}
\usepackage{amsmath,amssymb,mathtools}
\usepackage[scr=rsfs,cal=euler]{mathalfa}
\usepackage{xfrac}
\usepackage[unicode,hidelinks,bookmarks=false]{hyperref}
\newcommand{\ie}{i.e.\ }
\newcommand{\cf}{cf.\ }
\newcommand{\resp}{respectively\ }
\newcommand{\Subsection}{\subsection}
\providecommand{\nobreakdash}{\nobreak\hbox{-}}
\setlength{\emergencystretch}{2em}
\multlinegap0pt
\usepackage{bm}
\usepackage{bbm}
\usepackage{dsfont}
\usepackage{xspace}
\usepackage{cancel}
\usepackage{mathtools}
\usepackage{amsthm}
\makeatletter
\@ifundefined{theorem}{\newtheorem{theorem}{Theorem}}{}
\@ifundefined{corollary}{\newtheorem{corollary}[theorem]{Corollary}}{}
\makeatother
\DeclareMathOperator{\RM}{RM}
\newcommand{\CSS}{\textnormal{CSS}}
\newcommand{\F}{{\mathbb{F}}}
\newcommand{\ev}{\textnormal{ev}_{\F_2^m}}
\newcommand{\floor}[1]{\left\lfloor#1\right\rfloor}
\newcommand{\polyring}[1]{\F_2[x_1,\dots,x_m]^{#1}}
\newcommand{\punct}{\textnormal{punct}}
\newcommand{\short}{\textnormal{short}}
\newcommand{\supp}{\textnormal{supp}}
\title[CSS-T Codes from Reed Muller Codes]{CSS-T Codes from Reed Muller Codes}
\author{Emma Andrade, Jessalyn Bolkema, Thomas Dexter, Harrison Eggers, Victoria Luongo, Felice Manganiello, Luke Szramowski}
\date{Source: arXiv:2305.06423v2 (CC BY 4.0)}
\begin{document}
\maketitle

\noindent\emph{Source and licence.} CSS-T Codes from Reed Muller Codes. Version arXiv:2305.06423v2 (CC BY 4.0), distributed under the Creative Commons Attribution 4.0 International licence, as recorded in the arXiv licence metadata for this exact version and shown on the abstract page https://arxiv.org/abs/2305.06423v2. Full rights notice: licenses/arxiv-2305.06423-CC-BY-4.0.txt.\par\smallskip

\noindent\textit{Editorial scope.} Counted unit: the Theorem whose source label is \texttt{t:css-t} in the article (source0.tex lines 285--288, source label \texttt{t:css-t}), delivered together with the article's own complete original proof of it (source0.tex lines 290--320, one \texttt{proof} environment of 31 lines). No mathematics is written, repaired or re-derived: statement and proof are the article's own text line for line. Required context retained uncounted: cor:short (lines 270--273). Every definition, display and label of those lines is retained unchanged. Omitted from this package and not redistributed: 1--269, 274--284, 289--289, 321--414. Nothing inside a retained excerpt is omitted or altered: every retained body is delivered exactly as the article's lines are written, so no omission receipt of any kind accompanies this package. Citations: the retained excerpts carry no citation command at all, so no bibliography entry of the article is transcribed and no reference is redistributed. Reference adaptation: none was needed and none was made; every reference inside the delivered unit resolves inside it, so no unresolved reference or citation is produced. Printed slips kept literal: no printed word, symbol, hypothesis or number of the counted statement or of its proof was altered. Layout and class: the article's own class is \texttt{amsart} with packages including inputenc, amsmath, amsthm, amssymb, color, todonotes, stmaryrd; the wrapper is the lane's pinned \texttt{amsart} editorial wrapper at 10pt on A4, which restates the theorem-like environments the retained text needs and adds the article's own macro definitions that the retained text needs (\texttt{\textbackslash CSS}, \texttt{\textbackslash F}, \texttt{\textbackslash RM}, \texttt{\textbackslash ev}, \texttt{\textbackslash floor}, \texttt{\textbackslash polyring}, \texttt{\textbackslash punct}, \texttt{\textbackslash short}, \texttt{\textbackslash supp}), with extra wrapper packages (bm, bbm, dsfont, xspace, cancel, mathtools). The delivered numbering is the unit's own. Assets: no retained span contains a figure environment, a graphics-inclusion command, a PDF-page inclusion, a table or a TikZ picture; no image file is copied into this package, the package declares no asset dependency and no image file is redistributed with it. Licence: version arXiv:2305.06423v2 of this article is distributed under the Creative Commons Attribution 4.0 International licence, as recorded in the arXiv licence metadata for this exact version and shown on the abstract page https://arxiv.org/abs/2305.06423v2. A full rights notice is delivered at licenses/arxiv-2305.06423-CC-BY-4.0.txt. Characters: every retained excerpt is pure ASCII, so the delivered engine, XeLaTeX with its default Latin Modern text font, reports no missing character.
\begin{corollary}
Let $C_1$ and $C_2$ be two linear codes and $C=\CSS(C_1,C_2)$ a CSS code. $C$ is a CSS-T code if and only if for all $x\in C_2$ it holds that \[\punct(C_1,\F_2^m\setminus \supp(x))\subseteq \short(C_1^\perp,\F_2^m\setminus \supp(x)).\]
\label{cor:short}
\end{corollary}

\begin{theorem}\label{t:css-t}
Let $C_1=RM(\floor{\frac{m-1}{2}} - t, m)$ and  $C_2=RM(r_2,m)$. Then $\CSS(C_1,C_2)$ is a CSS-T code if and only if $r_2\leq 2t+1$ for $m$ even, or 
$r_2\leq 2t$ for $m$ odd.
\end{theorem}

\begin{proof}
Let $x\in C_2$. Then there exists a polynomial $p\in \polyring{r_2}$ such that $x=\ev(p)$. For simplicity in notation, let $r_1=\floor{\frac{m-1}{2}} - t$, so $C_1^\perp=\RM(m-r_1-1,m)$. Let $y\in C_1$, so there exists a polynomial $q\in \polyring{r_1}$ such that $y=\ev(q)$. 

Consider the polynomials $p$ and $1-p$. Because we are working over the binary field, it holds that $\supp(1-p)$ and $\supp(p)$ partition the set $\F_2^m$, meaning that $\supp(1-p)\cup \supp(p)=\F_2^m$ and $\supp(1-p)\cap \supp(p)=\emptyset$. This gives us a tool to split $\supp(q)$; specifically, given the polynomial decomposition $q=(1-p)q+pq$, it holds that $\supp((1-p)q)\subseteq \supp(1-p)$ and $\supp(pq)\subseteq \supp(p)$. Moreover, by the linearity of the evaluation map, it follows that $\ev(q)=\ev((1-p)q)+\ev(pq)$. 

Recall that $\supp(p)=\supp(x)$ whenever $x=\ev(p)$. Then
% \begin{align*}
%     \punct(y,\F_2^m\setminus \supp(x))&= \punct(\ev(q),\F_2^m\setminus \supp(x))\\
%     &= \punct(\ev((1-p)q)+\ev(pq),\F_2^m\setminus \supp(x))\\
%     &= \punct(\ev(pq),\F_2^m\setminus \supp(x))\\&=\short(\ev(pq),\F_2^m\setminus \supp(x))
% \end{align*}

\begin{align*}
    \punct(y,\F_2^m\setminus \supp(x))&= \punct(\ev(q),\F_2^m\setminus \supp(p))\\
    &= \punct(\ev((1-p)q)+\ev(pq),\supp(1-p))\\
    &= \punct(\ev(pq), \supp(1-p))\\&=\short(\ev(pq),\F_2^m\setminus \supp(p))
\end{align*}

where the last equality holds because $\supp(pq)\subseteq \supp(p)$.

To conclude, it is enough to prove that $pq\in\polyring{m-r_1-1}$, implying that $\ev(pq)\in C_1^\perp$. Indeed, 
\begin{align*}
    \deg(pq)&\leq r_1+r_2 = \floor{\frac{m-1}{2}} - t +r_2\\
    &=\begin{cases}
        h-1-t+2t+1, & m=2h \\ 
        h-t+2t, & m=2h+1
    \end{cases}\\
    &=h+t=m-r_1-1.
\end{align*}
Thus by Corollary~\ref{cor:short}, $\CSS(C_1,C_2)$ is indeed CSS-T.
\end{proof}

\end{document}
